\subsection{Commitment Was Never the Missing Piece}
\label{sec:12.1.2}
The OECD's recommendation. The European Commission's expert-group guidelines. UNESCO's. The two Partnerships on AI. The Hiroshima process. The safety-institute network. The Global Digital Compact. Eight overlapping voluntary bodies, cataloged in section~\ref{sec:10.8} beside the Council of Europe's 2024 framework convention — the only instrument in that landscape that would bind a party to anything, once it is in force, which at this writing it is not.

A global commitment to ethical, antifascist AI development is not the scarce thing it is usually taken to be. Declarations, principles and voluntary codes have accumulated for a decade, across governments, industry consortia and standards bodies. Anyone calling for the world to agree that AI should be developed responsibly is calling for something that has already happened, repeatedly, at the level of stated intent.

What has stayed scarce is the much narrower list of bodies that can act once a stated commitment is broken. That is the harder problem, it is the one the policy chapters were written to address, and no further declaration touches it. A commitment nobody can enforce is not a weak version of a commitment somebody can. It is a different kind of object, and counting the two together is what makes the landscape look fuller than it is. Section~\ref{sec:3.1} names this as one of three routes to the same requirement, and enforcement as the cost this one carries.

There is a third kind of object, and the argument of this book turns on it. A commitment can be declared, which is the decade cataloged above. It can be enforced from outside, which is what the policy chapters are for and what one instrument in that list even attempts. Or it can be held — carried as a reason by a party the thing at stake matters to, so that abandoning it is something that party would have to decide to do rather than something an absent enforcer failed to prevent. Chapter~\ref{sec:3} argues that a floor has to be the third kind, and this section supplies half of the reason from its own direction: the first kind binds nobody, and the second binds only where an enforcer exists, has jurisdiction, and is not itself the party in need of binding.

What it costs to build something that can hold a commitment is the whole of chapter~\ref{sec:3}, and the cost is a party that can be wronged. Set against that, the catalog above is not an argument for paying it. It is a measurement of what the alternative has delivered in ten years of trying, which is the comparison anyone weighing the cost should be making.
